\section{Results}
\label{sec:selected-observed-results}
% Generated from the frozen v4.9.12 result ledger; do not edit by hand.
\newcommand{\FinalSixteenSupportLowMeV}{30}
\newcommand{\FinalSixteenSupportHighMeV}{210}
\newcommand{\FinalSixteenLengthScaleUpperFactor}{12}
\newcommand{\SelectedAllThreeMinMass}{66}
\newcommand{\SelectedAllThreeMinP}{2.845\times 10^{-3}}
\newcommand{\SelectedAllThreeMinZ}{2.765}
\newcommand{\SelectedSharedEpsTwoHat}{4.550\times 10^{-6}}
\newcommand{\SelectedSharedEpsTwoSigma}{1.649\times 10^{-6}}


The results use the full 2015 and 2016 datasets and the available 10\%
sample from 2021.

Each mass point uses the corrected, bounded, piecewise-asymptotic \CLs{}
construction with the common signal-strength coordinate constrained to
$\eps^2\geq 0$. The Gaussian-process supports, hyperparameter states, and physics
normalizations are frozen before the observed limits and $p_0$ values are evaluated.
The inference is conditional on those frozen choices and on
the asymptotic approximation. For 2016, the frozen GP support is
\FinalSixteenSupportLowMeV--\FinalSixteenSupportHighMeV~MeV and the upper
resolution-scaled length-scale factor is \FinalSixteenLengthScaleUpperFactor. The
2015 and 2021 configurations retain their separately qualified supports.

The three datasets can be analyzed in seven distinct combinations: each dataset
alone, each of the three pairs, and all three together.

\begin{center}
\small
\begin{tabularx}{\linewidth}{@{}P{0.15\linewidth}Y P{0.28\linewidth}@{}}
\toprule
Scope & Active datasets & Statistical coordinate \\
\midrule
Standalone & full 2015 & campaign-specific $\eps^2$ \\
Standalone & full 2016 & campaign-specific $\eps^2$ \\
Standalone & 2021 10\% & campaign-specific $\eps^2$ \\
Pair & full 2015 + full 2016 & one shared $\eps^2$ \\
Pair & full 2015 + 2021 10\% & one shared $\eps^2$ \\
Pair & full 2016 + 2021 10\% & one shared $\eps^2$ \\
All three & full 2015 + full 2016 + 2021 10\% & one shared $\eps^2$ \\
\bottomrule
\end{tabularx}
\end{center}

The observed limits are shown first, followed by the pointwise expected
bands and empirical limit-tail probabilities. The background-only $p_0$ values are fixed-mass, one-sided asymptotic
quantities. The added Sidak curves give a separate analytic look-elsewhere
reference. Neither reference constitutes a direct coverage calibration of the
upper-limit construction.

\subsection{The complete observed search and conditional expected bands}
\label{sec:v5-band-results}

The full search runs from 19 to 250~MeV. At each mass it includes every campaign
whose declared search interval covers that point: 2015 alone at 19--38~MeV,
2015+2016 at 39--49~MeV, all three at 50--100~MeV, 2016+2021 at 101--180~MeV,
and 2021 alone at 181--250~MeV. This membership rule is fixed by the search
ranges; it never selects the dataset giving the smallest observed limit.
Figure~\ref{fig:v5-union-bands} places the nominal local asymptotic $p_0$ directly
below the corresponding observed limit and conditional expected bands, followed
by an aligned significance panel. Black gives the local reference and red its
Sidak correction for this complete search.

The bands use 300 background-only draws per mass from the
mass-local GP prediction. Each draw first samples a correlated latent
background and then Poisson counts; constituent draws are shared among the
combinations at that mass. The GP state is held fixed. Different masses do not
form one coherent full-spectrum experiment. The resulting quantiles describe
pointwise variation under that conditional model; they do not calibrate the
\CLs{} sampling distributions or supply a global significance.

\fig{0.98}{../figures/v5_total_limits_with_local_pvalues.pdf}{Complete 19--250~MeV
observed search, with the preassigned active datasets shown above the curves.
The upper panel shows 90\% \CLs{} limits and the median, central 68\% and central
95\% ranges of the 300 conditional background-only toys per mass, with newly computed
bands at the added 2015 coordinates.
The next two panels show the fixed-mass asymptotic $p_0$ and
its normal-quantile significance, with the local reference in black and the
Sidak reference in red. The latter uses $N_{\rm eff}=35.381$ for the
complete search; its value at 66~MeV is $Z_{\rm Sidak}=1.305$,
compared with local $Z=2.765$. The observed
polyline connects every adjacent tested mass, including membership transitions;
no numerical endpoint is changed by those connections. The physical
visible-branching correction is applied once above the dimuon threshold.
These are pointwise bands; 2016-inclusive values retain the numerical exception.}
{fig:v5-union-bands}

\paragraph{Reading the significance panels.}
We use the one-sided normal convention $Z=\Phi^{-1}(1-p)$ and display
$Z^+=\max(0,Z)$, with integer guides on the vertical axis. Thus a probability
of at least one half appears at zero; this convention avoids interpreting a
negative normal quantile as positive discovery evidence. It does not change
the probability. The signed fit coordinate $r$, which can describe a deficit,
is shown separately in Section~\ref{sec:v5-global-results}.

The Sidak curves apply Eq.~\eqref{eq:sidak-global} to each local threshold
using one fixed $N_{\rm eff}$ for the entire corresponding search.
The inherited resolution-spacing prescription sums
$\Delta m/(2.25\,\bar\sigma_m)$ between adjacent mass hypotheses, uses the
narrowest active campaign resolution, and caps the result between one and
the number of tested masses. It gives 14.521, 16.062 and 26.402 for the
individual 2015, 2016 and 2021 searches. This is an analytic reference for
the look-elsewhere penalty, not a calibration from complete simulated scans.
Table~\ref{tab:v502-sidak-lookup} shows how the same local significance maps
to different corrections as the search range and resolution change.

\begin{table}[htbp]\centering\small
\begin{tabular}{rrrrrr}\toprule
Local $Z$ & Local $p$ & 2015 & 2016 & 2021 & Complete search\\
 & & \multicolumn{4}{c}{Sidak $Z^+$}\\\midrule
1 & 0.159 & 0.000 & 0.000 & 0.000 & 0.000\\
2 & 0.0228 & 0.571 & 0.499 & 0.112 & 0.000\\
3 & 0.00135 & 2.066 & 2.024 & 1.811 & 1.678\\
4 & 3.17e-05 & 3.314 & 3.286 & 3.143 & 3.056\\
5 & 2.87e-07 & 4.457 & 4.435 & 4.327 & 4.262\\
\bottomrule\end{tabular}

\caption{Local-to-Sidak significance conversion for the declared individual
and complete searches. These are illustrative thresholds passed through the
fixed analytic formula, not additional observed peaks. The final four columns
use the $Z^+$ display convention; signed quantiles and probabilities are
retained in the accompanying numerical tables.}
\label{tab:v502-sidak-lookup}
\end{table}

\fig{0.98}{../figures/v5_observed_coupling_overlay.pdf}{Observed coupling limits
from the three individual released samples and the full maximal-available-dataset
combination of Fig.~\ref{fig:v5-union-bands}. Only observed limits are shown.
The combination is obtained from a shared-coupling likelihood with the fixed
membership rule; it is not the envelope of the individual curves.}
{fig:v5-observed-overlay}

\FloatBarrier

\paragraph{The 2015 90--100 MeV extension.}
All 2015-containing curves now include the ten added integer hypotheses.
The 14--135~MeV training support is retained and the 90~MeV kernel parameters
are held fixed above that endpoint. The template width follows the existing
linear response, rising from 4.70 to 5.23~MeV across the added interval.
The shaded range has new observed limits, 300 conditional toys per mass and
scope, and the saved 95~MeV recovery check. Appendix~\ref{app:v504-extension}
records the calculation and its validation limits. Historical displays are
retained in Appendix~\ref{app:v504-historical-displays}.


\subsection{Conditional full-2021 response compared with BaBar}
\label{sec:v501-babar-proxy}
To show how the current result responds to a larger 2021 sample, we repeat
the restricted density-response construction used in Figure~84 of v4.9.7.
Let $\rho_d(m)$ be the observed prompt normalization density in the physical
window $m\pm1.64\sigma_m$. Holding the full 2015 and 2016 samples fixed and
increasing only the 2021 contribution by a factor of ten gives
\begin{equation}
 (\epsilon^2_{90})_{\rm proxy}(m)
 = (\epsilon^2_{90})_{\rm current}(m)
 \sqrt{\frac{\sum_{d\in\mathcal D(m)}\rho_d(m)}
 {\sum_{d\in\mathcal D(m)}s_d\rho_d(m)}} ,
 \qquad (s_{2015},s_{2016},s_{2021})=(1,1,10).
 \label{eq:v501-density-proxy}
\end{equation}
Here $\mathcal D(m)$ is the set of contributing datasets. The scale is unity below
50~MeV, density weighted where samples overlap, and exactly $1/\sqrt{10}$
where only 2021 contributes. The current and proxy minima are
$8.0043\times10^{-7}$ and $3.1898\times10^{-7}$, both at 72~MeV.
These are properties of the displayed curves, not expected future outcomes.

\fig{0.98}{../figures/v501_current_projected_babar.pdf}{Current observed
90\% upper limits and the conditional full-2021 density-response proxy
compared with the published BaBar visible-dark-photon observed 90\% contour
\cite{BaBarDarkPhoton2014}. The gray dashed curve is exactly the current
maximal-available-dataset HPS result over 19--250~MeV. The blue curve applies
Eq.~\eqref{eq:v501-density-proxy} to its current observed normalization
densities. It retains present fluctuations and is neither expected
sensitivity nor a future observed limit or a fit to full-2021 data.
HPS and BaBar use different statistical constructions. The coupling conversion includes the dimuon branching correction.}
{fig:v501-babar-proxy}
\FloatBarrier

\subsection{Signal-yield and coupling limits for each dataset}
\label{sec:v501-individual-scans}\label{sec:selected-individual-results}
Figures~\ref{fig:v501-scan-2015}--\ref{fig:v501-scan-2021} separate the
observed yield limit from its conversion to $\epsilon^2$. The upper panels
use the full Gaussian-template yield, while the likelihood uses the portion
inside the fitted window. The different mass dependence of the yield and
coupling limits follows from the dataset's mass-dependent normalization.
The nominal local minima occur at 51, 90, and 78~MeV in 2015, 2016, and
2021, with $p_0=8.46\times10^{-4}$, $3.09\times10^{-4}$, and
$2.47\times10^{-3}$, respectively. They are fixed-mass asymptotic quantities;
the coherent global calculation in Section~\ref{v5-sec:global-method} uses
a separately declared scan statistic and conditional background model.

\fig{0.95}{../figures/v501_observed_scan_2015.pdf}{Observed 2015 full-sample
scan on 19--100~MeV with 14--135~MeV GP support. From top to bottom: the
90\% \CLs{} full-template signal-yield upper limit $A_{90}$, the corresponding
$\epsilon^2$ upper limit, discovery probabilities, and aligned $Z^+$ values.
The local asymptotic reference is black and the Sidak reference is red,
with $N_{\rm eff}=14.521$. At 51~MeV they give $Z=3.139$ and
$Z_{\rm Sidak}=2.250$. The added 91--100~MeV fits use the fixed endpoint prescription; the
shaded interval has the conditional status in Appendix~\ref{app:v504-extension}.
The Sidak reference is distinct from the conditional GP global calculation
in Section~\ref{sec:v5-global-results}.}
{fig:v501-scan-2015}
\fig{0.95}{../figures/v501_observed_scan_2016.pdf}{Corresponding observed
2016 full-sample scan on 39--180~MeV with 30--210~MeV GP support. The panels
show the full-template $A_{90}$, $\epsilon^2$ limit, discovery probabilities,
and aligned $Z^+$ values. Black is local and red is the Sidak reference
with $N_{\rm eff}=16.062$; their values at 90~MeV are 3.424 and 2.579.
}
{fig:v501-scan-2016}
\fig{0.95}{../figures/v501_observed_scan_2021.pdf}{Corresponding observed
2021 10\% scan on 50--250~MeV with 36--300~MeV GP support. The panels
show the full-template $A_{90}$, minimal-visible $\epsilon^2$ limit,
discovery probabilities, and aligned $Z^+$ values. Black is local and red
is the Sidak reference with $N_{\rm eff}=26.402$; their values at 78~MeV
are 2.810 and 1.527. Above the dimuon threshold the inherited
branching correction enters the coupling conversion once; it does not change
the visible electron-pair yield or local probability.}
{fig:v501-scan-2021}

\begin{table}[htbp]\centering\small
\begin{tabular}{lrrrrr}\toprule
Search & Mass [MeV] & $N_{\rm eff}$ & Local $Z$ & Sidak $p$ & Sidak $Z$\\\midrule
All three & 66 & 9.948 & 2.765 & 0.02794 & 1.912\\
2015 & 51 & 14.521 & 3.139 & 0.01222 & 2.250\\
2016 & 90 & 16.062 & 3.424 & 0.00495 & 2.579\\
2021 (10\%) & 78 & 26.402 & 2.810 & 0.06333 & 1.527\\
2015 + 2016 & 90 & 10.011 & 3.659 & 0.00127 & 3.019\\
2015 + 2021 & 50 & 9.948 & 3.351 & 0.00400 & 2.652\\
2016 + 2021 & 66 & 20.708 & 2.484 & 0.12627 & 1.144\\
Complete search & 66 & 35.381 & 2.765 & 0.09589 & 1.305\\
\bottomrule\end{tabular}

\caption{Largest local discovery significance in each declared search and
the corresponding Sidak reference. Each row uses its own fixed search range
and effective trials count. The complete search uses the preassigned
maximal-available-dataset rule. Selecting the most favorable row would add
another search and is not included in these corrections. These values use
the nominal asymptotic results, rather than the separately evaluated
stress-centered scores.}
\label{tab:v502-sidak-peaks}
\end{table}

\FloatBarrier
\subsection{Combining datasets through a shared coupling}
\label{sec:selected-combined-results}

Figure~\ref{fig:combined-final-results} contains all three pairwise combinations and
the all-three result. At any given mass, the simultaneous likelihood profiles one
nonnegative $\eps^2$ shared by the active campaigns. Each campaign nevertheless
retains its own signal normalization, GP prediction, covariance, and observed
spectrum. Combining the datasets is thus a new likelihood calculation, not an
envelope, average, or choice among the standalone limits. Each curve is evaluated only
on the mass interval common to all contributing datasets.

\begin{figure}[H]
  \centering
  \includegraphics[width=0.94\linewidth,height=0.59\textheight,keepaspectratio]{../figures/combined_final_results.pdf}
  \caption{Observed shared-$\eps^2$ results for the three pairwise combinations
  and the all-three overlap. For visual alignment, the top coordinate expresses the
  shared-coupling limit as the total compatible signal yield summed over the active
  campaigns. The middle panel gives the 90\% bounded, piecewise-asymptotic \CLs{}
  upper limit on $\eps^2$, and the bottom panel gives the local asymptotic $p_0$.
  The 2015-containing overlaps extend through 100~MeV. The shaded interval
marks the extension described in Appendix~\ref{app:v504-extension}.
No expected-limit bands are shown.}
  \label{fig:combined-final-results}
\end{figure}

The shared coupling expresses the physical relationship among the campaigns: a
putative dark photon has one $\eps^2$, while its expected event count
changes with beam conditions, exposure, acceptance, and mass resolution. It also
keeps the independent background information from each spectrum explicit.

\FloatBarrier
\subsection{Local asymptotic scans}
\label{sec:selected-local-p0-series}

Figure~\ref{fig:final-asymptotic-pvalues} places the local $p_0$ scans for all seven
cases on the same axes, making structure at a common mass easier to compare across
campaigns. Every plotted point is still a separate fixed-mass test. The conditional GP global studies are discussed separately in
Section~\ref{sec:v5-global}; they do not change these local asymptotic curves.
The separate black-local/red-Sidak panels for all four overlapping
combinations appear in Figure~\ref{fig:v502-combination-discovery},
with their peak conversions in Table~\ref{tab:v502-sidak-peaks}.

\begin{figure}[htbp]
  \centering
  \includegraphics[width=0.98\linewidth]{../figures/final_asymptotic_pvalues.pdf}
  \caption{Local one-sided asymptotic background-only $p_0$ for the three
  individual datasets, the three pairs, and the all-three combination. The values are
  conditional on the frozen GP supports and states. No look-elsewhere correction,
  global significance, or toy calibration is supplied by these curves.
The 2015-containing scopes include 91--100~MeV.}
  \label{fig:final-asymptotic-pvalues}
\end{figure}

\FloatBarrier
The common three-campaign interval has its smallest local asymptotic value at
\SelectedAllThreeMinMass~MeV, with $p_0=\SelectedAllThreeMinP$ and
$Z=\SelectedAllThreeMinZ$. The current detailed extraction displays appear in
Section~\ref{sec:v5-extractions}.
\FloatBarrier
